%%%%%%%%%%%%%%%%%%%%%%%%%%%%%%%%%%%%%%%%%%%%%%%%%%%%%%%%%%%%%%
%%%%%%%%%%%% CA-0411 Análisis de Datos I %%%%%%%%%%%%%%%%%%%%%%
%%%%%%%%%%%%%%%%%%%%%%%%%%%%%%%%%%%%%%%%%%%%%%%%%%%%%%%%%%%%%%
\newpage
\subsection*{CA-0411 Análisis de Datos I}
\addcontentsline{toc}{subsection}{CA-0411 Análisis de Datos I}

\hypertarget{CA0411}{}
\begin{refsection}
\begin{table}[H]
\centering{
\begin{tabular}{|ll|}
\hline
\textbf{Año:} IV  & \textbf{Requisitos:} CA-0307 \\
\textbf{Ciclo:} I  & \textbf{Correquisitos:} Ninguno \\
\textbf{Tipo de curso:} Teórico-práctico & \textbf{Horas semanales:} 5 \\
\textbf{Créditos:} 4 & \textbf{Virtualidad:} PV, V \\ \hline
\end{tabular}
}
\end{table}

\subsubsection*{Descripción del curso}

Este curso está dirigido a estudiantes avanzados que han completado los dos cursos de estadística actuarial por lo que se ubica en el cuarto año y primer ciclo. El curso pretende introducir al estudiante a la modelación de aprendizaje estadístico y la implementación de modelos de aprendizaje no supervisado y supervisado y que el/la estudiante aprenda herramientas de modelación para análisis de datos que tienen aplicaciones en múltiples disciplinas incluyendo riesgos financieros, seguros y finanzas.

\subsubsection*{Objetivo general}

Proporcionar al estudiante la teoría detrás del aprendizaje estadístico para la implementación y análisis de modelos de aprendizaje supervisados y no supervisados para datos reales.

\subsubsection*{Objetivos específicos}

Al finalizar el curso el/la estudiante estará en la capacidad de:

\begin{enumerate}
\item Comprender el modelo de aprendizaje estadístico y su relación con el proceso de modelación en Ciencia de Datos. (SC11, SH05.)
\item Implementar modelos de aprendizaje no supervisado y supervisado en conjuntos de datos de distinta índole. (SC11, SH05.)
\item Analizar los procesos de construcción de modelos de aprendizaje no supervisado y supervisado con especial énfasis en aplicaciones actuariales, financieras y demográficas. (SC11, SH05.)
\end{enumerate}

\subsubsection*{Contenidos}

\begin{enumerate}
\item \textbf{Introducción a aprendizaje estadístico.}
\item \textbf{Aprendizaje no supervisado de datos}
\begin{enumerate}
\item Reducción en dimensión a través de análisis en componentes principales.
\item Clustering o análisis de conglomerados: k-medias y clasificación jerárquica.
\end{enumerate}
\item \textbf{Introducción al aprendizaje supervisado}
\begin{enumerate}
\item Análisis discriminante.
\item Máquinas de soporte vectorial.
\item Árboles de decisión.
\end{enumerate}
\end{enumerate}

\subsubsection*{Metodología}

Se espera que la persona docente aplique al menos tres de las siguientes estrategias didácticas:

\begin{enumerate}
\item Proyecto.
\item Avances de proyecto (bitácoras).
\item Búsqueda de datos y preguntas de investigación en fuentes externas.
\item UVE heurística de Gowin.
\item Foros y ubicación de conceptos en contexto.
\item Elaboración de un wiki para identificación de conceptos relevantes.
\item Visitas de expertos.
\item Clases magistrales: sincrónicas y asincrónicas. Peer project learning.
\item Identificación de conceptos en artículos científicos sencillos e interpretación de resultados.
\item Laboratorios guiados en R o Python.
\item Solución de problemas reales ligados a la industria.
\end{enumerate}

Se motiva a la persona docente a que incentive la investigación con un alto nivel computacional en este curso a través de un proyecto que involucre trabajo de los/las estudiantes durante todas las 15 semanas del curso. Además se debe incentivar la aplicación de modelos estadísticos en la resolución de problemas que involucren la cuantificación de riesgos de distinta índole, y finalmente la interpretación de resultados hacia público no especializado.

\subsubsection*{Evaluación}

\begin{enumerate}
\item Exámenes.
\item Documentos sintéticos: póster, resumen (\emph{abstract}), etc.
\item Presentaciones orales.
\item Rúbricas.
\item Ensayos cortos de noticias, gráficos.
\item Desarrollo de software: paquetes o módulos de R o Python.
\end{enumerate}

\subsubsection*{Bibliografía}

\begin{enumerate}
\item James, G., Witten, D., Hastie, T., y Tibshirani, R. (2013). \emph{An Introduction to Statistical Learning with Applications in R.} Springer.
\item Trejos, J.; Castillo, W.; González, J. (2014). \emph{Análisis Multivariado de Datos. Métodos y Aplicaciones.} Editorial de la Universidad de Costa Rica, San José.
\item Hastie, T., Tibshirani, R., y Friedman, J. (2009). \emph{The Elements of Statistical Learning: Data Mining, Inference, and Prediction} (2nd ed.). Springer.
\item Hair, J. F., Black, W. C., Babin, B. J., y Anderson, R. E. (2019). \emph{Multivariate Data Analysis} (8th ed.). Cengage.
\item Lebart, L.; Morineau, A.; y Piron, M. (1995). \emph{Statistique Exploratoire Multidimensionnelle.} Dunod, Paris.
\item Escofier, B., y Pagès, J. (1998). \emph{Analyses factorielles simples et multiples.} Dunod, Paris.
\item Pagès, J. (2015). \emph{Multiple Factor Analysis by Example Using R.} Chapman and Hall/CRC.
\item Chatfield, C., y Collins, A. J. (1980). \emph{Introduction to Multivariate Analysis.} Chapman and Hall (incluye escalamiento multidimensional).
\end{enumerate}

\end{refsection}
